% GENERATED by tools/gen.py from the unit's vocabulary -- do not edit.
Ćwiczenia tego bloku są \emph{generowane} z listy słów tygodnia i sprawdzają pamięć słownictwa, nie gramatykę. Jeśli odpowiedź jest synonimem, a nie słowem z listy, sprawdź w tabeli powyżej.
\begin{cwiczeniadalej}
\exitem{Podaj rodzajnik: \dots{} \de{Hilfe} (pomoc)}{\artf{die} Hilfe}
\exitem{Podaj rodzajnik: \dots{} \de{Pause} (przerwa)}{\artf{die} Pause}
\exitem{Podaj rodzajnik: \dots{} \de{Regel} (reguła)}{\artf{die} Regel}
\exitem{Podaj rodzajnik: \dots{} \de{Verbot} (zakaz)}{\artn{das} Verbot}
\exitem{Podaj rodzajnik: \dots{} \de{Pflicht} (obowiązek)}{\artf{die} Pflicht}
\exitem{Podaj rodzajnik: \dots{} \de{Aufgabe} (zadanie)}{\artf{die} Aufgabe}
\exitem{Podaj rodzajnik: \dots{} \de{Kunde} (klient)}{\artm{der} Kunde}
\exitem{Podaj rodzajnik: \dots{} \de{Meeting} (spotkanie (robocze))}{\artn{das} Meeting}
\exitem{Podaj rodzajnik: \dots{} \de{Parkplatz} (miejsce parkingowe)}{\artm{der} Parkplatz}
\exitem{Podaj liczbę mnogą: \de{\artf{die} Pause}}{\de{\artf{die} Pausen}}
\exitem{Podaj liczbę mnogą: \de{\artf{die} Regel}}{\de{\artf{die} Regeln}}
\exitem{Podaj liczbę mnogą: \de{\artn{das} Verbot}}{\de{\artf{die} Verbote}}
\exitem{Podaj liczbę mnogą: \de{\artf{die} Pflicht}}{\de{\artf{die} Pflichten}}
\exitem{Podaj liczbę mnogą: \de{\artf{die} Aufgabe}}{\de{\artf{die} Aufgaben}}
\exitem{Podaj liczbę mnogą: \de{\artm{der} Kunde}}{\de{\artf{die} Kunden}}
\exitem{Podaj liczbę mnogą: \de{\artn{das} Meeting}}{\de{\artf{die} Meetings}}
\exitem{Podaj liczbę mnogą: \de{\artm{der} Parkplatz}}{\de{\artf{die} Parkplätze}}
\exitem{Jak to jest po niemiecku: musieć?}{\de{müssen}}
\exitem{Jak to jest po niemiecku: mieć powinność, mieć zrobić?}{\de{sollen}}
\exitem{Jak to jest po niemiecku: pomagać?}{\de{helfen}}
\exitem{Jak to jest po niemiecku: pokazywać?}{\de{zeigen}}
\exitem{Jak to jest po niemiecku: wysyłać?}{\de{schicken}}
\exitem{Jak to jest po niemiecku: przerwa?}{\de{\artf{die} Pause, Pausen}}
\exitem{Jak to jest po niemiecku: zakaz?}{\de{\artn{das} Verbot, Verbote}}
\exitem{Jak to jest po niemiecku: klient?}{\de{\artm{der} Kunde, Kunden}}
\exitem{Jak to jest po niemiecku: niestety?}{\de{leider}}
\exitem{Jak to jest po niemiecku: natychmiast?}{\de{sofort}}
\exitem{Co znaczy \de{können}?}{móc; umieć}
\exitem{Co znaczy \de{dürfen}?}{wolno}
\exitem{Co znaczy \de{helfen}?}{pomagać}
\exitem{Co znaczy \de{erklären}?}{wyjaśniać}
\exitem{Co znaczy \de{zeigen}?}{pokazywać}
\exitem{Co znaczy \de{\artf{die} Hilfe}?}{pomoc}
\exitem{Co znaczy \de{\artf{die} Pause, Pausen}?}{przerwa}
\exitem{Co znaczy \de{\artn{das} Verbot, Verbote}?}{zakaz}
\exitem{Co znaczy \de{\artf{die} Aufgabe, Aufgaben}?}{zadanie}
\exitem{Co znaczy \de{\artm{der} Kunde, Kunden}?}{klient}
\exitem{Podaj formę \de{du} czasownika \de{können}}{\de{du kannst}}
\exitem{Podaj formę \de{du} czasownika \de{müssen}}{\de{du musst}}
\exitem{Podaj formę \de{du} czasownika \de{wollen}}{\de{du willst}}
\exitem{Podaj formę \de{du} czasownika \de{dürfen}}{\de{du darfst}}
\exitem{Podaj formę \de{du} czasownika \de{sollen}}{\de{du sollst}}
\end{cwiczeniadalej}
